% No preamble. Finite-regulator quantum-domain interface; qg- labels.
\section{Size of the flat constrained Gaussian fluctuations}
\label{qg-section}

There is a quantitative reason why the preceding smooth-background
ball cannot itself be used as a uniform domain for integrating all
fine Gaussian fluctuations. We prove it already at the flat
reference, in four dimensions. Let $n_f=M^4n_c$ be the numbers
of fine and coarse sites of a compatible periodic regulator,
where $M=L^j\ge L\ge3$ and $\eta=M^{-1}$. Work in one fixed
real trace-orthonormal color component; the flat operator acts
identically on every color. Set
\[
 X=\ker Q\cap\ker(RD^T),\qquad K=d^Td.
\]
The positive restricted Maxwell form defines a centered real
Gaussian $Z\in X$ with density proportional to
$\exp[-\langle Z,KZ\rangle_\eta/2]$ relative to the
physical Euclidean measure on $X$. Its covariance operator
under the physical pairing is the constrained inverse
$\mathcal C$ used above.

\begin{lemma}\label{qg-dimension}
For each real color component,
\[
 \dim X=3(n_f-n_c),\qquad
 \operatorname{Tr}(d\mathcal C d^T)=\dim X.
\]
\end{lemma}
\begin{proof}
The map $Q$ on bonds is onto because $QH=I$, hence
$\dim\ker Q=4(n_f-n_c)$. The scalar projection $R$ has
rank $n_f-n_c$: its complementary projection has the full
coarse scalar rank $n_c$. The map $RD^T$ from $\ker Q$
onto $\operatorname{ran}R$ is surjective. Indeed the scalar
operator $T=G'R$ satisfies $Q_0T=0$ and $R\mathcal L T=R$,
where $\mathcal L=D^TD$. For $f\in\operatorname{ran}R$,
the bond field $DTf$ has $QDTf=D_cQ_0Tf=0$ and
$RD^TDTf=f$. Rank subtraction gives the dimension.

Let $J_X$ embed $X$ isometrically into the physical bond
Hilbert space, and let $K_X=J_X^TKJ_X$. Then
$\mathcal C=J_XK_X^{-1}J_X^T$. Cyclicity of the finite
operator trace gives
$\operatorname{Tr}(d\mathcal C d^T)
=\operatorname{Tr}(K_XK_X^{-1})=\dim X$.
\end{proof}

The physical covariance identity is
$\mathbb E\langle f,Z\rangle_\eta\langle g,Z\rangle_\eta
=\langle f,\mathcal Cg\rangle_\eta$.
Thus the covariance of the actual coordinate values is
$\eta^{-4}\mathcal C$, and that of the physical curls is
$\eta^{-4}d\mathcal C d^T$. There are $6n_f$ positive
plaquette components per color. Lemma~\ref{qg-dimension}
therefore gives the exact average
\begin{equation}\label{qg-curl-variance}
 \boxed{\quad
 \frac1{6n_f}\sum_p\mathbb E|(dZ)_p|^2
     =\frac{M^4-1}{2}.
 \quad}
\end{equation}
In particular at least one real curl coordinate has variance
at least $(M^4-1)/2$, irrespective of the periodic volume.
This conclusion uses the physical covariance normalization;
omitting $\eta^{-4}$ would give a different and incorrect scale.

Fix a deterministic background $A_0$, a radius $r>0$ and
a coupling factor $g>0$. The event
$N_\alpha(A_0+gZ)\le r$ bounds every physical first difference,
and hence bounds the selected curl coordinate by $C r$,
where the color-norm constant is fixed. It restricts a
one-dimensional centered Gaussian of standard deviation
at least $c\eta^{-2}$ to an interval of length at most
$C r/g$; the interval center may depend on $A_0$.
The maximum of the Gaussian density therefore proves
\begin{equation}\label{qg-small-ball}
 \boxed{\quad
 \mathbb P\{N_\alpha(A_0+gZ)\le r\}
       \le \min\left(1,\frac{C r\eta^2}{g}\right).
 \quad}
\end{equation}
The bound is uniform in the deterministic background,
refinement and compatible periodic volume. Thus for any
sequence $g=g_\eta$ with $g_\eta/\eta^2\to\infty$, a fixed
physical smooth ball has probability tending to zero under
this constrained flat Gaussian reference.

This is an obstruction to integrating all fluctuations inside
the minimizer's fixed $N_\alpha$ ball, not an obstruction to
a Yang--Mills theory. Equation~\eqref{qg-curl-variance} shows
that the dimensionless microscopic curl $\eta^2dZ$ has
bounded averaged variance. A quantum renormalization step
must instead control suitable rescaled fluctuation fields
and their tails while retaining smooth estimates for the
background. Averaged or component variance bounds alone
do not control a supremum over a growing number of variables,
nor do they bound interacting polymer activities. No
non-Gaussian probability estimate is inferred here.
